%Daten zur Institution

%\input{Dozentdaten}

%\renewcommand{\fachbereich}{Fachbereich}

%\renewcommand{\dozent}{Prof. Dr. . }

%Klausurdaten

\renewcommand{\klausurgebiet}{ }

\renewcommand{\klausurtyp}{ }

\renewcommand{\klausurdatum}{. 20}

\klausurvorspann {\fachbereich} {\klausurdatum} {\dozent} {\klausurgebiet} {\klausurtyp}

%Daten für folgende Punktetabelle

\renewcommand{\aeins}{  3  }

\renewcommand{\azwei}{  3   }

\renewcommand{\adrei}{  5  }

\renewcommand{\avier}{  2  }

\renewcommand{\afuenf}{  4  }

\renewcommand{\asechs}{  8  }

\renewcommand{\asieben}{  4  }

\renewcommand{\aacht}{  5  }

\renewcommand{\aneun}{  3  }

\renewcommand{\azehn}{  2  }

\renewcommand{\aelf}{  6  }

\renewcommand{\azwoelf}{  7  }

\renewcommand{\adreizehn}{  4   }

\renewcommand{\avierzehn}{  8  }

\renewcommand{\afuenfzehn}{  64  }

\renewcommand{\asechzehn}{    }

\renewcommand{\asiebzehn}{    }

\renewcommand{\aachtzehn}{    }

\renewcommand{\aneunzehn}{    }

\renewcommand{\azwanzig}{    }

\renewcommand{\aeinundzwanzig}{    }

\renewcommand{\azweiundzwanzig}{    }

\renewcommand{\adreiundzwanzig}{    }

\renewcommand{\avierundzwanzig}{    }

\renewcommand{\afuenfundzwanzig}{    }

\renewcommand{\asechsundzwanzig}{    }

\punktetabellevierzehn

\klausurnote

\newpage

\setcounter{section}{0}

\inputaufgabegibtloesung
{3}
{

Definiere die folgenden
\zusatzklammer {kursiv gedruckten} {} {} Begriffe.
\aufzaehlungsechs{Die
\stichwort {zweite tangentiale Ableitung} {}
\zusatzklammer {oder die
\stichwort {tangentiale Beschleunigung} {}} {} {}
einer
\zusatzklammer {als Abbildung nach $\R^n$} {} {}
zweimal
\definitionsverweis {stetig differenzierbaren Kurve}{}{}
\maabbdisp {\gamma} {I} {Y
} {}
in eine
\definitionsverweis {differenzierbare Hyperfläche}{}{}

\mavergleichskette
{\vergleichskette
{ Y
}
{ \subseteq }{ \R^n
}
{  }{
}
{    }{
}
{   }{
}
}
{}{}{.}

}{Ein \stichwort {zusammenhängender} {} topologischer Raum $X$.

}{Die \stichwort {Tangentialabbildung in einem Punkt} {}

\mavergleichskette
{\vergleichskette
{  P
}
{ \in }{ M
}
{  }{}
{    }{}
{   }{}
}
{}{}{}
zu einer differenzierbaren Abbildung
\maabbdisp {\varphi} { M } { N
} {,}
wobei
\mathkor {} {M} {und} {N} {}
differenzierbare Mannigfaltigkeiten sind.

}{Das
\stichwort {Produkt} {}
von differenzierbaren Mannigfaltigkeiten
\mathkor {} {M} {und} {N} {.}

}{Das
\stichwort {Volumenmaß} {}
zu einer positiven Volumenform $\omega$ auf einer $n$-dimensionalen differenzierbaren Mannigfaltigkeit.

}{Eine
\stichwort {geodätische Kurve} {}
auf einer riemannschen Mannigfaltigkeit $M$.
}

}
{} {}

\inputaufgabegibtloesung
{3}
{

Formuliere die folgenden Sätze.
\aufzaehlungdrei{Der
\stichwort {Satz über die Krümmung auf einen implizit gegebenen ebenen Kurve} {}

\mavergleichskette
{\vergleichskette
{  Y
}
{ \subseteq }{  \R^2
}
{  }{
}
{    }{
}
{   }{
}
}
{}{}{.}}{Die
\stichwort {Transformationsformel für positive Volumenformen auf Mannigfaltigkeiten} {.}}{Der Satz über die Gaußkrümmung und die Schnittkrümmung auf einer differenzierbaren Fläche

\mavergleichskette
{\vergleichskette
{  Y
}
{ \subseteq }{  \R^3
}
{  }{
}
{    }{
}
{   }{
}
}
{}{}{}}

}
{} {}

\inputaufgabegibtloesung
{5}
{

Es sei

\mavergleichskette
{\vergleichskette
{ G
}
{ \subseteq }{ \R^n
}
{  }{
}
{    }{
}
{   }{
}
}
{}{}{}
\definitionsverweis {offen}{}{}
und
\maabbdisp {\varphi} { G } { \R^m
} {}
eine
\definitionsverweis {stetig differenzierbare Abbildung}{}{,}
die im Punkt

\mavergleichskette
{\vergleichskette
{ P
}
{ \in }{ G
}
{  }{
}
{    }{
}
{   }{
}
}
{}{}{}
ein surjektives
\definitionsverweis {totales Differential}{}{}
besitze. Es sei

\mavergleichskette
{\vergleichskette
{ v
}
{ \in }{ T_PZ
}
{  }{
}
{    }{
}
{   }{
}
}
{}{}{}
ein Vektor des
\definitionsverweis {Tangentialraumes an die Faser}{}{}
$Z$ zu $\varphi$ durch $P$. Zeige, dass es eine
\definitionsverweis {stetig differenzierbare Kurve}{}{}
\maabbdisp {\gamma} { ]- \epsilon, \epsilon[ } { Z
} {}
\zusatzklammer {für ein geeignetes

\mavergleichskettek
{\vergleichskettek
{ \epsilon
}
{ > }{ 0
}
{  }{
}
{    }{
}
{   }{
}
}
{}{}{}} {} {}
mit

\mavergleichskette
{\vergleichskette
{ \gamma(0)
}
{ = }{ P
}
{  }{
}
{    }{
}
{   }{
}
}
{}{}{}
und mit

\mavergleichskettedisp
{\vergleichskette
{ \gamma'(0)
}
{ =} { v
}
{ } {
}
{ } {
}
{ } {
}
}
{}{}{}
gibt.

}
{} {}

\inputaufgabegibtloesung
{2}
{

Es seien
\mathkor {} {r} {und} {R} {}
positive reelle Zahlen mit

\mavergleichskette
{\vergleichskette
{ 0
}
{ < }{ r
}
{ < }{ R
}
{    }{
}
{   }{
}
}
{}{}{.}
Bestimme ein
\definitionsverweis {Einheitsnormalenfeld}{}{}
für den
\zusatzklammer {eingebetteten} {} {}
Torus

\mavergleichskettedisp
{\vergleichskette
{ T
}
{ =} { \left\{ (x,y,z) \in \R^3  \mid   \left(  \sqrt{x^2+y^2} -R  \right)^2 +z^2  = r^2         \right\}
}
{ } {
}
{ } {
}
{ } {
}
}
{}{}{.}

}
{} {}

\inputaufgabe
{4}
{

Man erläutere die Relevanz des Satzes über implizite Abbildungen für den Aufbau der Theorie der Mannigfaltigkeiten.

}
{} {}

\inputaufgabegibtloesung
{8}
{

Beweise den Satz über die Selbstadjungiertheit der Weingartenabbildung.

}
{} {}

\inputaufgabegibtloesung
{4}
{

Es sei $X$ ein
\definitionsverweis {kompakter Raum}{}{}
und es sei

\mavergleichskette
{\vergleichskette
{Y
}
{ \subseteq }{ X
}
{  }{
}
{    }{
}
{   }{
}
}
{}{}{}
eine
\definitionsverweis {abgeschlossene Teilmenge}{}{,}
die die
\definitionsverweis {induzierte Topologie}{}{}
trage. Zeige, dass $Y$ ebenfalls kompakt ist.

}
{} {}

\inputaufgabegibtloesung
{5}
{

Man gebe ein Beispiel einer
\definitionsverweis {zweidimensionalen}{}{}
\definitionsverweis {zusammenhängenden}{}{} \definitionsverweis {differenzierbaren Mannigfaltigkeit}{}{}
$M$ und einem Punkt

\mavergleichskette
{\vergleichskette
{ P
}
{ \in }{ M
}
{  }{
}
{    }{
}
{   }{
}
}
{}{}{}
derart, dass
\mathkor {} {M} {und} {M \setminus \{P\}} {}
zueinander
\definitionsverweis {diffeomorph}{}{}
sind.

}
{} {}

\inputaufgabegibtloesung
{3}
{

Man gebe für jeden
\definitionsverweis {Tangentialvektor}{}{}

\mavergleichskette
{\vergleichskette
{ v
}
{ \in }{ T_PS^2
}
{ \subseteq }{ \R^3
}
{    }{
}
{   }{
}
}
{}{}{}
mit

\mavergleichskette
{\vergleichskette
{ \Vert {v} \Vert
}
{ = }{ 1
}
{  }{
}
{    }{
}
{   }{
}
}
{}{}{}
in einem Punkt

\mavergleichskette
{\vergleichskette
{ P
}
{ \in }{ S^2
}
{ \subset }{ \R^3
}
{    }{
}
{   }{
}
}
{}{}{}
auf der
\definitionsverweis {Einheitssphäre}{}{}
einen
\definitionsverweis {differenzierbaren}{}{}
Repräsentanten
\maabbdisp {\gamma} { I } { S^2
} {}
mit

\mavergleichskette
{\vergleichskette
{ \gamma(0)
}
{ = }{ P
}
{  }{
}
{    }{
}
{   }{
}
}
{}{}{}
an.

}
{} {}

\inputaufgabegibtloesung
{2}
{

Wir betrachten die Standardparabel

\mavergleichskette
{\vergleichskette
{ Y
}
{ \subseteq }{ \R^2
}
{  }{
}
{    }{
}
{   }{
}
}
{}{}{}
mit der induzierten riemannschen Struktur und mit der Parametrisierung
\maabbeledisp {\varphi} {\R} { Y
} { t } { \left(  t   , \,  t^2                   \right)
} {,}
die wir als eine
\zusatzklammer {inverse} {} {}
Karte betrachten. Bestimme die riemannsche
\definitionsverweis {Fundamentalfunktion}{}{}

\mavergleichskette
{\vergleichskette
{ g(t)
}
{ = }{ g_{11} (t)
}
{  }{
}
{    }{
}
{   }{
}
}
{}{}{.}

}
{} {}

\inputaufgabegibtloesung
{6}
{

Berechne die zurückgezogene Differentialform
\mathl{\varphi^* \tau}{} zu

\mavergleichskettedisp
{\vergleichskette
{ \tau
}
{ =} { dx \wedge dy \wedge dz - w dx \wedge dy \wedge dw + \cos (xy)  dx \wedge dz \wedge dw-yw dy \wedge dz \wedge dw
}
{ } {
}
{ } {
}
{ } {
}
}
{}{}{}
unter der Abbildung
\maabbeledisp {\varphi} { \R^3 } { \R^4
} { (r,s,t) } {  \left(  r^2s,t, \sin  r ,e^{st}  \right)  = (x,y,z,w)
} {.}

}
{} {}

\inputaufgabegibtloesung
{7 (2+3+2)}
{

Wir betrachten die differenzierbaren Abbildungen
\maabbeledisp {\gamma} { [1,2] } {\R^2
} { t } {(t,t^{-1})
} {,}
und
\maabbeledisp {\varphi} {\R^2} {\R^3
} {(u,v)} {(u^2,uv,-u+v^2)
} {,}
und die Differentialform

\mavergleichskettedisp
{\vergleichskette
{ \omega
}
{ =} { xdx-zdy+dz
}
{ } {
}
{ } {
}
{ } {
}
}
{}{}{}
auf dem $\R^3$.

a) Berechne die zurückgezogene Differentialform
\mathl{\varphi^* \omega}{} auf dem $\R^2$.

b) Berechne das Wegintegral zur Differentialform
\mathl{\varphi^* \omega}{} zum Weg $\gamma$.

c) Berechne
\zusatzklammer {ohne Bezug auf b)} {} {}
das Wegintegral zur Differentialform
\mathl{\omega}{} zum Weg $\varphi \circ \gamma$.

}
{} {}

\inputaufgabegibtloesung
{4}
{

Es sei $M$ eine
\definitionsverweis {Mannigfaltigkeit mit Rand}{}{.}
Zeige, dass jede offene Teilmenge

\mavergleichskette
{\vergleichskette
{ N
}
{ \subseteq }{ M
}
{  }{
}
{    }{
}
{   }{
}
}
{}{}{}
ebenfalls eine Mannigfaltigkeit mit
\zusatzklammer {eventuell leerem} {} {}
Rand ist, und dass

\mavergleichskettedisp
{\vergleichskette
{ \partial N
}
{ =} { \partial M \cap N
}
{ } {
}
{ } {
}
{ } {
}
}
{}{}{}
gilt.

}
{} {}

\inputaufgabegibtloesung
{8}
{

Beweise den Satz über die Charakterisierung von geodätischen Kurven bezüglich des Levi-Civita-Zusammenhangs.

}
{} {}

 [[Kategorie:Latexseite]]
